\section{研究背景及问题重述}
\subsection{研究背景}
神经网络推理通常由大量带有数据依赖的矩阵、向量和搬运操作组成。将完整计算图放在单个核心上执行可以减少跨核通信，但计算单元和片上存储的并行能力没有被充分利用；将计算图切分到多个核心又会引入边界搬运、同步等待和共享 DDR 争用。因此，多核调度的关键不是单独增加并行度，而是在并行收益、数据驻留和搬运代价之间作出可检查的取舍。

题目给出的计算图采用 Tensor--Op 二部结构。输入和最终输出张量位于 DDR，中间张量位于每个核心私有的 L1 或 UB；计算操作占用 Cube 或 Vector 流水线，COPY 操作占用搬运流水线。评估程序根据参赛者提交的切图和调度方案自动补充边界 COPY、进行核内容量处理，并在共享带宽下模拟所有核心的执行。因此，论文中的方案必须与评估程序的对象、时序和容量语义一致。

本文使用“驻留域”表示一份数据可以直接复用的计算范围。问题一的驻留域是一个独立子图，问题二和问题三的驻留域是一个核心上的合并 Task。这个定义使三问都可以先回答“一个张量要进入多少个驻留域”，再分别处理 Task 同步、私有容量和共享 Cache 查询时序。
\resultbox{18mm}{三问硬件场景、驻留域与数据流关系图（待绘制方法图）}
\begin{table}[H]\centering\caption{题面对象与本文建模对象的对应关系}
\begin{tabularx}{\linewidth}{lYY}\toprule
题面对象 & 本文记号或状态 & 在算法中的作用\\\midrule
计算图 & $G=(T,V,E)$ & 给出张量、操作和依赖关系\\
子图 & $s$ & 切图和 Task 构造的基本单元\\
核心 & $k\in\{0,\ldots,K-1\}$ & 承载子图并形成核内顺序\\
张量 & $t$，大小 $b_t$ & 决定驻留、搬运和容量占用\\
评估器 & $F_r(X)$ & 对场景 $r$ 的官方 Makespan 进行统一评价\\\bottomrule
\end{tabularx}\end{table}

\subsection{问题重述}
问题一要求在场景 A 中为所有非 COPY 操作分配子图 id，并为每个核心给出子图执行顺序。每个子图独立形成一个 Task，跨子图数据必须通过 DDR 中转。本文需要输出可执行的切图方案、2 至 5 核的调度方案以及与单核基准同口径的平均加速比和逐用例指标；实验结果暂不在本稿填入。

问题二要求在场景 B 中将同一核心上的子图合并为一个 Task，使同核数据可以保留在 L1/UB 中，同时对跨核数据插入 COPY\_OUT/COPY\_IN 和同步延迟。本文在问题一的切图基础上，增加容量峰值和张量存活区间的约束，并通过核内顺序调整降低不必要的换入换出。

问题三在场景 B 的 Task 结构上增加所有核心共享的只读 L2 Cache。题面给定 Cache 容量和带宽，要求分析其资源作用并设计对应的切图与调度算法。本文把共享输入的消费者聚合、消费者分散和查询时序调整作为三类候选来源，并分别保留无 L2 与有 Cache 的同口径比较位置。

\subsection{全文组织}
第 2 节给出三问输入、输出和统一计算链；第 3 节定义符号和评价量；第 4 节说明切图、分核和顺序之间的难点；第 5 节列出必要假设及失效条件；第 6、7 节分别建立问题一/二和问题三的模型；第 8 节说明三问共用的候选生成、官方展开和离散事件模拟；第 9 节给出待补结果的组织方式、稳健性检查和局限；后文列出来源、支撑材料和 AI 使用记录。
